\documentclass[11pt,a4paper]{article}
\usepackage[utf8]{inputenc}
\usepackage{amsmath,amssymb,amsthm}
\usepackage{algorithm}
\usepackage{algpseudocode}
\usepackage{booktabs}
\usepackage{hyperref}
\usepackage{geometry}
\geometry{margin=1in}

\title{Temporal Convolutional Networks (TCN): Dilated Causal Convolutions, Exponential Receptive Fields, and Gated WaveNet Dynamics}
\author{Automorphic Intelligence Architecture Specifications}
\date{\today}

\newtheorem{theorem}{Theorem}
\newtheorem{lemma}[theorem]{Lemma}
\newtheorem{proposition}[theorem]{Proposition}
\theoremstyle{definition}
\newtheorem{definition}{Definition}
\newtheorem{remark}{Remark}

\begin{document}

\maketitle

\begin{abstract}
Recurrent sequence models suffer from sequential computational bottlenecks during backpropagation through time and struggle with long-horizon gradient degradation. Temporal Convolutional Networks (TCNs) combine 1D causal convolutions (preventing future information leakage) with exponentially increasing dilation factors ($d = 2^l$), endowing pure feed-forward convolutional stacks with exponential receptive field growth. Augmented with gated activation functions and residual identity mappings, TCNs enable massive GPU-parallel sequence training while maintaining streaming autoregressive capability. This document formalizes the causal dilated discrete convolution operator, proves exact exponential receptive field scaling bounds, details algorithmic pseudocode, and delivers a verified step-by-step numerical calculation.
\end{abstract}

\section{Notation and Mathematical Preliminaries}

Let $\mathbf{x} = (\mathbf{x}_1, \dots, \mathbf{x}_T)$ denote an input sequence where $\mathbf{x}_t \in \mathbb{R}^C$. A 1D convolutional filter is parameterized by kernel size $K \ge 2$, dilation rate $d \ge 1$, and weights $\mathbf{W} \in \mathbb{R}^{C_{\text{out}} \times C_{\text{in}} \times K}$.

\begin{table}[h]
\centering
\caption{Mathematical Notation for Temporal Convolutional Networks}
\begin{tabular}{lll}
\toprule
\textbf{Symbol} & \textbf{Domain} & \textbf{Semantic Description} \\
\midrule
$T$ & $\mathbb{N}$ & Temporal sequence length \\
$K$ & $\mathbb{N}_{\ge 2}$ & 1D discrete spatial/temporal kernel width (typically $K=2$ or $K=3$) \\
$l \in \{0, \dots, L-1\}$ & $\mathbb{N}$ & Layer depth index in the stacked architecture \\
$d_l = 2^l$ & $\mathbb{N}$ & Exponential dilation rate at layer $l$ \\
$*_{d}$ & Operator & Dilated discrete convolution operator with stride 1 \\
$\mathcal{R}_L$ & $\mathbb{N}$ & Total cumulative receptive field span after $L$ layers \\
$\sigma(\cdot), \tanh(\cdot)$ & Activation & WaveNet gated non-linear activation functions \\
\bottomrule
\end{tabular}
\end{table}

\section{Problem Definition and Core Concepts}

\subsection{3-Level Explanation}
\begin{enumerate}
    \item \textbf{Intuitive (High School level):} An RNN listens one second at a time, unable to fast-forward. A TCN looks at the whole audio recording at once with binoculars that double their zoom at each layer: layer 1 sees 2 steps back, layer 2 sees 4, layer 3 sees 8, layer 4 sees 16—covering hours of sound in a split second without peeking into the future.
    \item \textbf{Engineering Level:} TCN applies 1D convolutions with causal left-padding of $(K-1)d$ zeros so time $t$ never sees time $t+1$. Doubling dilation rates $d \in \{1, 2, 4, 8, 16\}$ yields exponential history coverage with $\mathcal{O}(L)$ depth instead of $\mathcal{O}(T)$ sequential unrolling.
    \item \textbf{Mathematical Level:} The causal dilated convolution operator is defined by:
    \begin{equation}
    \mathbf{y}_t = (\mathbf{W} *_{d} \mathbf{x})_t = \sum_{k=0}^{K-1} \mathbf{W}[k] \cdot \mathbf{x}_{t - d \cdot k}
    \end{equation}
    combined with WaveNet gated activations:
    \begin{equation}
    \mathbf{z}_t = \tanh(\mathbf{W}_f *_{d} \mathbf{x})_t \odot \sigma(\mathbf{W}_g *_{d} \mathbf{x})_t
    \end{equation}
\end{enumerate}

\subsection{5-Step Algorithmic Framework}
\begin{enumerate}
    \item \textbf{Causal Left-Padding:} Pad input tensor on the left with $(K - 1)d$ zeros to enforce causality and maintain length $T$.
    \item \textbf{Dilated Strided Gathering:} Sample inputs at dilated intervals $t, t-d, t-2d, \dots$.
    \item \textbf{Filter and Gate Projection:} Compute parallel filter activations $\mathbf{f}_t = \tanh(\text{Conv}_d)$ and gate activations $\mathbf{g}_t = \sigma(\text{Conv}_d)$.
    \item \textbf{Hadamard Elementwise Gating:} Modulate filter by gate: $\mathbf{z}_t = \mathbf{f}_t \odot \mathbf{g}_t$.
    \item \textbf{Residual Addition and Stacking:} Project $\mathbf{z}_t$ via $1 \times 1$ conv and add identity shortcut: $\mathbf{y}_t = \mathbf{x}_t + W_{\text{res}} \mathbf{z}_t$.
\end{enumerate}

\section{Algorithm Specification}

\begin{algorithm}[H]
\caption{TCN Causal Dilated Residual Forward Pass}
\label{alg:tcn}
\begin{algorithmic}[1]
\Require Sequence $\mathbf{X} \in \mathbb{R}^{T \times C}$, kernel size $K$, layer depths $L$, weights $\{W_f^{(l)}, W_g^{(l)}, W_{\text{res}}^{(l)}\}_{l=0}^{L-1}$.
\Ensure Output sequence $\mathbf{Y} \in \mathbb{R}^{T \times C}$.
\State $\mathbf{H}^{(0)} \gets \mathbf{X}$
\For{$l = 0$ \textbf{to} $L-1$}
    \State $d \gets 2^l$
    \State $\mathbf{F} \gets \tanh(\text{CausalDilatedConv1D}(\mathbf{H}^{(l)}, W_f^{(l)}, d))$
    \State $\mathbf{G} \gets \sigma(\text{CausalDilatedConv1D}(\mathbf{H}^{(l)}, W_g^{(l)}, d))$
    \State $\mathbf{Z} \gets \mathbf{F} \odot \mathbf{G}$
    \State $\mathbf{H}^{(l+1)} \gets \mathbf{H}^{(l)} + \text{Linear1x1}(\mathbf{Z}, W_{\text{res}}^{(l)})$
\EndFor
\State \Return $\mathbf{H}^{(L)}$
\end{algorithmic}
\end{algorithm}

\section{Theoretical Guarantees and Receptive Field Scaling}

\begin{theorem}[Exponential Receptive Field Growth]
\label{thm:tcn_rf}
For an $L$-layer TCN with kernel size $K$ and exponentially doubling dilation rates $d_l = 2^l$, the total cumulative receptive field $\mathcal{R}_L$ satisfies:
\begin{equation}
\mathcal{R}_L = 1 + (K - 1) \sum_{l=0}^{L-1} 2^l = 1 + (K - 1)(2^L - 1) = \mathcal{O}(K \cdot 2^L)
\end{equation}
guaranteeing exponential history coverage with linear parameter depth $\mathcal{O}(L)$.
\end{theorem}

\begin{proof}
At layer 0 ($d_0 = 1$), the receptive field is $r_0 = K$. For layer $l$, each of the $K$ kernel points spans a distance $(K - 1) 2^l$ relative to the layer below. By induction:
\begin{equation}
\mathcal{R}_L = 1 + \sum_{l=0}^{L-1} (K - 1) 2^l = 1 + (K - 1) \sum_{l=0}^{L-1} 2^l
\end{equation}
Applying the geometric series formula $\sum_{l=0}^{L-1} 2^l = 2^L - 1$ yields $\mathcal{R}_L = 1 + (K - 1)(2^L - 1)$.
\end{proof}

\section{Complexity Analysis}

\begin{itemize}
    \item \textbf{Time Complexity:} $\mathcal{O}(T \cdot L \cdot K \cdot C^2)$ operations (fully parallel across sequence length $T$ during training).
    \item \textbf{Space Complexity:} $\mathcal{O}(T \cdot L \cdot C)$ intermediate activation storage for full parallel backpropagation.
\end{itemize}

\section{Exactness and Error Analysis}

Because convolutions operate on structured 1D tensors with static left-padding, causality is mathematically guaranteed ($\frac{\partial y_t}{\partial x_{t + \tau}} \equiv 0$ for all $\tau > 0$), preventing any validation leakage.

\section{Worked Numerical Example}

Let $K=2, L=3, C=1$.
Dilation rates: $d_0 = 1, d_1 = 2, d_2 = 4$.
Receptive field: $\mathcal{R}_3 = 1 + (2 - 1)(2^3 - 1) = 1 + 1(7) = 8$ steps.
Consider sequence $X = [1.0, 2.0, 3.0, 4.0]$ ($T=4$).
Let filter kernel $W = [0.5, 0.5]$ at layer 0 ($d=1$):
\begin{itemize}
    \item $t=0$: $x_0=1.0$, padded $x_{-1}=0.0 \implies y_0 = 0.5(1.0) + 0.5(0.0) = 0.5$.
    \item $t=1$: $x_1=2.0, x_0=1.0 \implies y_1 = 0.5(2.0) + 0.5(1.0) = 1.5$.
    \item $t=2$: $x_2=3.0, x_1=2.0 \implies y_2 = 0.5(3.0) + 0.5(2.0) = 2.5$.
    \item $t=3$: $x_3=4.0, x_2=3.0 \implies y_3 = 0.5(4.0) + 0.5(3.0) = 3.5$.
\end{itemize}

\section{Numerical Considerations and Edge Cases}

\begin{itemize}
    \item \textbf{Autoregressive Fast Generation via Circular Queues:} While training is parallel, sequential sample-by-sample inference can cache the past $(K-1)d$ activations in a FIFO ring buffer, reducing per-step inference latency to $\mathcal{O}(L \cdot K \cdot C^2)$ instead of recomputing past convolutions.
\end{itemize}

\section{References}
\begin{enumerate}
    \item van den Oord, A., Dieleman, S., Zen, H., Simonyan, K., Vinyals, O., Graves, A., ... & Kavukcuoglu, K. (2016). WaveNet: A generative model for raw audio. \emph{SSW}, 125.
    \item Bai, S., Kolter, J. Z., & Koltun, V. (2018). An empirical evaluation of generic convolutional and recurrent networks for sequence modeling. \emph{arXiv preprint arXiv:1803.01271}.
\end{enumerate}

\end{document}
